\section{Conclusion}
\label{sec:chapter8-conclusion}

This thesis presented Stockrium, an integrated decision-support and experimentation platform for the Vietnamese equity market. The project addresses a practical information problem: an individual investor must often combine price behavior, company fundamentals, and financial news from separate sources before forming a judgment. Stockrium organizes these evidence domains through two complementary clients, a shared backend, structured data-acquisition pipelines, and a coordinated multi-agent workflow. The mobile application supports frequent market exploration and individual stock analysis, while Stockrium Lab provides a larger web workspace for longer-running AI analyses and historical backtests.

The principal contribution is the implemented user workflow rather than a claim of investment performance. An authenticated Stockrium Lab user can configure a supported run, monitor its status, retain the resulting record, reopen its details, and compare compatible saved results. Experiment configurations, summaries, detailed outputs, available reproducibility metadata, and artifact references remain associated with their owner. The system therefore lets users inspect what was executed and how outputs differ instead of exposing only an isolated recommendation or return value.

The backend separates presentation, domain services, persistence, deterministic financial computation, and language-model reasoning. Technical, fundamental, and news agents examine different evidence domains, after which deterministic controls combine their scores into a Buy-or-Wait assessment and the aggregator produces a structured explanation. This separation makes the contributing evidence visible and allows the same analytical services to support both interactive mobile use and Stockrium Lab. Stockrium remains a decision-support tool: it neither executes securities orders nor replaces the user's judgment or professional financial advice.

The representative VN30 case in Chapter~6 demonstrates that the web workflow can produce, preserve, inspect, and compare historical outputs. Its mixed results, including cases in which simple references performed better, are retained as part of the record. These values illustrate the transparency and operation of the tool under documented assumptions; they do not prove superior returns, reliable market outperformance, or future profitability.

\section{Achievement of the Project Objectives}
\label{sec:objective-achievement}

The three objectives introduced in Section~\ref{sec:application-objectives} were achieved within the functional and data boundaries defined for the project. Table~\ref{tab:objective-achievement} relates each objective to the corresponding implemented outcome.

\begin{table}[H]
    \centering
    \caption{Achievement of the Stockrium project objectives}
    \label{tab:objective-achievement}
    \begin{tabularx}{\textwidth}{p{1.85cm} p{3.2cm} X}
        \toprule
        \textbf{Objective} & \textbf{Assessment} & \textbf{Implemented outcome} \\
        \midrule
        O1 & Achieved within scope & The mobile application and shared backend integrate market discovery, stock evidence, configurable multi-agent analysis, explainable Buy-or-Wait assessments, watchlists, and conversational assistance into an end-to-end workflow for individual investors. \\
        O2 & Achieved within scope & Stockrium Lab allows authenticated users to run supported AI analyses and backtests, follow their status, retain user-owned records, reopen detailed results, and compare two to five compatible records. The workflow includes ownership enforcement, warnings, and available reproducibility metadata. \\
        O3 & Achieved within scope & Acquisition and normalization pipelines maintain the deployed Vietnamese-market dataset covering OHLCV observations, company and financial information, industry data, and financial articles with processed summaries and sentiment labels. \\
        \bottomrule
    \end{tabularx}
\end{table}

Objective O1 is supported by a coherent path from market discovery to evidence inspection and AI-assisted synthesis. Users may rely on horizon-based automatic settings or manually select analytical perspectives and their relative weights. The result separates the assessment and proposed trading-plan fields from the technical, fundamental, and news explanations, helping the user review the basis of the output rather than treating it as an opaque instruction.

Objective O2 is supported by the implemented Stockrium Lab lifecycle. AI analyses and backtests create persisted records with a type, scope, configuration, status, result, timestamps, and available version information. Long-running backtests return control to the web client while their status is polled, and history and comparison operations remain restricted to the owning account. The representative case confirms that a user can complete this workflow; historical returns are outputs displayed by the tool, not criteria used to declare the objective successful.

Objective O3 is supported by the data pipelines used by both clients and the historical simulator. The deployed data are normalized into structures required by the technical, fundamental, and news services and are refreshed through maintainable collection processes. Source availability, storage limits, publication timing, and live-data changes still constrain completeness and exact reproducibility; these boundaries are stated as limitations rather than hidden behind a general claim of data accuracy.

\section{Future Work}
\label{sec:future-work}

\subsection{Stockrium Lab Reliability and Reproducibility}

The current asynchronous backtest implementation persists status but relies on in-process background execution. A server restart can interrupt an active run, and cancellation, durable resume, automatic retry, and queued resource control are not yet complete. Future work should move long-running jobs to a durable worker queue, preserve symbol-level checkpoints, expose cancellation and retry controls, and provide clear recovery states in the history interface.

Saved metadata identifies the configuration and available application, code, pipeline, prompt, model, and data-snapshot versions, but live sources and model-backed outputs cannot always be reproduced exactly. Stronger reproducibility requires immutable input snapshots, versioned prompts and models, explicit random or sampling controls where supported, and retention policies that keep each result linked to the evidence used at run time. The comparison view can then distinguish an exact rerun from a configuration-only reconstruction.

Stockrium Lab can also be extended with clearer naming, tagging, search, and grouping of saved runs; richer configuration differences; and exportable comparison reports. These additions should preserve the current rule that only records of the same type are compared and that every history, detail, comparison, and artifact operation is scoped to the authenticated owner.

\subsection{Analytical Grounding and Human Review}

The current explanations do not attach every claim to an exact candlestick, reporting period, ratio, or article. Future specialist-agent outputs should include structured evidence references that the clients can open directly. The aggregation stage should validate each material claim against those references, identify missing or conflicting evidence, and report uncertainty without converting a weakly supported result into a confident recommendation.

Broader human review is also needed. The professional feedback reported in this thesis provides one useful perspective, but it is not a controlled usability study or a broad assessment by investment experts. Future studies should ask retail investors and financial professionals to complete defined mobile and Stockrium Lab tasks, then assess factual grounding, explanation clarity, comparison usefulness, error recognition, and the degree to which the tool supports further investigation. Such work should evaluate usability and decision-support quality without using short-term profitability as a proxy for user value.

\subsection{Historical Simulation and Data Coverage}

The current simulator assigns capital independently to each security and simplifies gaps, price-limit queues, liquidity, bid--ask spreads, slippage, partial fills, and market impact. Future work should add a portfolio-level cash balance, position sizing, concurrent-exposure limits, portfolio risk controls, and more realistic fill assumptions. These changes would make the displayed historical records more representative of the conditions a user expects to inspect, while still not turning a backtest into a promise of live performance.

The representative case covers the current VN30 stock set over a limited interval. Future demonstrations can include longer periods, point-in-time index membership, more varied market regimes, and broader supported universes where data quality permits. Exact publication timestamps for financial statements and immutable news snapshots would reduce look-ahead and reconstruction uncertainty. Broader cases should remain demonstrations of tool behavior unless a separate study supplies the statistical design and evidence required for a market-performance claim.

\subsection{System Quality and User Experience}

The repository currently has substantial automated backend coverage but no automated React Native component, mobile end-to-end, or browser suite. Future work should add tests for critical mobile navigation and analysis flows, Stockrium Lab authentication and saved-result workflows, accessibility, responsive layouts, and failure states. Security assessment, load and stress testing, dependency-failure exercises, and recovery testing are also necessary before production-scale deployment.

The application should become more resilient to market-data, storage, and language-model service failures through caching, fallback providers, health monitoring, stale-data indicators, and graceful degradation. Real-time subscriptions may be introduced for appropriate mobile market views with explicit reconnection and freshness handling. These improvements should prioritize reliable access to existing decision-support functions rather than expanding Stockrium into brokerage execution or automatic trading.

\section{Closing Remarks}
\label{sec:closing-remarks}

Stockrium demonstrates an end-to-end environment in which Vietnamese-market information, multi-agent analysis, and historical outputs can be accessed through interfaces suited to different user tasks. Its mobile application supports routine exploration, and Stockrium Lab allows users to run, retain, reopen, and compare more demanding analyses on the web. The resulting transparency---including saved configurations, supporting evidence, unfavorable outcomes, and stated limitations---is the central result of the project.

Accordingly, the thesis concludes that Stockrium provides a practical foundation for inspectable financial decision support and user-driven experimentation. It does not conclude that the analytical pipeline is consistently profitable. Future progress should be judged by stronger evidence traceability, reproducibility, reliability, usability, and honest comparison, while final investment responsibility remains with the user.
